\documentclass[10pt,a4paper]{amsart}
\usepackage[margin=19mm]{geometry}
\usepackage{amsmath,amssymb,mathtools}
\usepackage[scr=rsfs,cal=euler]{mathalfa}
\usepackage{xfrac}
\usepackage[unicode,hidelinks,bookmarks=false]{hyperref}
\newcommand{\ie}{i.e.\ }
\newcommand{\cf}{cf.\ }
\newcommand{\resp}{respectively\ }
\newcommand{\Subsection}{\subsection}
\providecommand{\nobreakdash}{\nobreak\hbox{-}}
\setlength{\emergencystretch}{2em}
\multlinegap0pt
\usepackage{bm}
\usepackage{bbm}
\usepackage{dsfont}
\usepackage{xspace}
\usepackage{cancel}
\usepackage{mathtools}
\usepackage{amsthm}
\makeatletter
\@ifundefined{subsection}{\newtheorem{subsection}{Subsection}}{}
\@ifundefined{lemma}{\newtheorem{lemma}[subsection]{Lemma}}{}
\makeatother
\title[Huber means on Riemannian manifolds]{Huber means on Riemannian manifolds}
\author{Jongmin Lee, Sungkyu Jung}
\date{Source: arXiv:2407.15764v2 (CC BY 4.0)}
\begin{document}
\maketitle

\noindent\emph{Source and licence.} Huber means on Riemannian manifolds. Version arXiv:2407.15764v2 (CC BY 4.0), distributed under the Creative Commons Attribution 4.0 International licence, as recorded in the arXiv licence metadata for this exact version and shown on the abstract page https://arxiv.org/abs/2407.15764v2. Full rights notice: licenses/arxiv-2407.15764-CC-BY-4.0.txt.\par\smallskip

\noindent\textit{Editorial scope.} Counted unit: the Lemma whose source label is \texttt{lem:subseq} in the article (source0.tex lines 1353--1355, source label \texttt{lem:subseq}), delivered together with the article's own complete original proof of it (source0.tex lines 1356--1376, one \texttt{proof} environment of 21 lines). No mathematics is written, repaired or re-derived: statement and proof are the article's own text line for line. Required context retained uncounted: none. Every definition, display and label of those lines is retained unchanged. Omitted from this package and not redistributed: 1--1352, 1377--3205. Nothing inside a retained excerpt is omitted or altered: every retained body is delivered exactly as the article's lines are written, so no omission receipt of any kind accompanies this package. Citations: the retained excerpts carry no citation command at all, so no bibliography entry of the article is transcribed and no reference is redistributed. Reference adaptation: none was needed and none was made; every reference inside the delivered unit resolves inside it, so no unresolved reference or citation is produced. Printed slips kept literal: no printed word, symbol, hypothesis or number of the counted statement or of its proof was altered. Layout and class: the article's own class is \texttt{imsart} with packages including fix-cm, latexsym, array, color, rotating, epsfig, wrapfig, graphicx, lscape, graphics, verbatim, fancyvrb, pdfpages, fmtcount; the wrapper is the lane's pinned \texttt{amsart} editorial wrapper at 10pt on A4, which restates the theorem-like environments the retained text needs and adds the article's own macro definitions that the retained text needs (none), with extra wrapper packages (bm, bbm, dsfont, xspace, cancel, mathtools). The delivered numbering is the unit's own. Assets: no retained span contains a figure environment, a graphics-inclusion command, a PDF-page inclusion, a table or a TikZ picture; no image file is copied into this package, the package declares no asset dependency and no image file is redistributed with it. Licence: version arXiv:2407.15764v2 of this article is distributed under the Creative Commons Attribution 4.0 International licence, as recorded in the arXiv licence metadata for this exact version and shown on the abstract page https://arxiv.org/abs/2407.15764v2. A full rights notice is delivered at licenses/arxiv-2407.15764-CC-BY-4.0.txt. Characters: every retained excerpt is pure ASCII, so the delivered engine, XeLaTeX with its default Latin Modern text font, reports no missing character.
\begin{lemma}\label{lem:subseq}
For a point $p \in M$, $p \in \cap_{n=1}^{\infty}\overline{\cup_{k=n}^{\infty}E_{k}}$ if and only if there exist an increasing sequence $k_{n} ~ (k_{n} \ge n)$ and a sequence of points $p_{n} \in E_{k_{n}}$ such that $p_{n} \underset{n \rightarrow \infty}{\rightarrow} p$ (i.e., $d(p_{n}, p) \underset{n \rightarrow \infty}{\rightarrow} 0$).
\end{lemma}

\begin{proof}[Proof of Lemma~\ref{lem:subseq}]
($\Rightarrow$) Denote by $A'$ the set of accumulation points of $A \subseteq M$. It is known that $\overline{A}=A \cup A'$ for any set $A \subset M$. Define $A_{n}=\cup_{k=n}^{\infty}E_k$. Then, $p \in \cap_{n=1}^{\infty}\overline{A_{n}}$ implies $p \in \overline{A_{n}}=A_{n} \cup A'_{n}$ for all $n\ge 1$, thereby meaning that
\[
p \in A_{n} ~ \mbox{infinitely many} ~ n \quad \mbox{or} \quad p \in A'_{n} ~ \mbox{infinitely many}~ n.
\]
When the former case occurs, let $p_{n}=p$ for all $n \ge 1$. When the later case occurs, suppose that $p \in A'_{n_{k}}$ for any $k \ge 1$. By the definition of accumulation point,
\begin{eqnarray*}
\exists p_{1}, k_{1}~(\ge n_1) \quad &\mbox{s.t.}& \quad p_1 \in E_{k_1},~ d(p_1,p)<1, \\
\exists p_{2}, k_{2}~(\ge n_2, k_1) \quad &\mbox{s.t.}& \quad p_2 \in E_{k_2}, ~ d(p_2,p)<\frac{1}{2}, \\
\exists p_{3}, k_{3}~(\ge n_3, k_2) \quad &\mbox{s.t.}& \quad p_3 \in E_{k_3}, ~ d(p_3,p)<\frac{1}{3}, \\
&\vdots& \\
\exists p_{\ell}, k_{\ell}~(\ge n_\ell, k_{\ell-1}) \quad &\mbox{s.t.}& \quad p_\ell \in E_{k_\ell}, ~ d(p_\ell,p)<\frac{1}{\ell}.
\end{eqnarray*}
Replacing the $\ell$ by $n$, we obtain that $p_{n} \in E_{k_{n}}$ for any $n \ge 1$, and that $d(p_{n}, p) \underset{n \rightarrow \infty}{\rightarrow} 0$.  

($\Leftarrow$) If $p_{n}=p$ infinitely many $n$, then $p \in A_{m} \subset \overline{A_{m}}$ for any integer $m \ge 1$, meaning that $p \in \cap_{m=1}^{\infty} \overline{A_{m}}$. If $p_{n} \neq p$ all but finitely many $n$, then $p \in A'_{m} \subseteq \overline{A_{m}}$ for any integer $m\ge 1$, meaning that $p \in \cap_{m=1}^{\infty} \overline{A_{m}}$. 



Combining these facts, we complete the proof.    
\end{proof}

\end{document}
